\documentclass[fullpage]{uiucthesis}

\usepackage[T1]{fontenc}

\begin{document}

\title{Learning Sparse Representations for Low-Resource\\
       Sequence Labelling}
\author{Priya Anand}
\department{Computer Science}
\schools{B.S., Purdue University, 2018\\
         M.S., University of Illinois at Urbana-Champaign, 2021}
\phdthesis
\advisor{Robert T. Halvorsen}
\degreeyear{2026}
\maketitle

\frontmatter

\begin{abstract}
This thesis studies sparse representation learning for sequence
labelling tasks in low-resource languages, where labelled training
data is scarce and standard dense embeddings overfit quickly. We
introduce a regularization scheme that encourages structured
sparsity in the hidden layers of a recurrent encoder, and show on
five typologically diverse languages that the resulting models
generalize substantially better than their dense counterparts when
fewer than one thousand labelled sentences are available. We further
analyze the learned sparsity patterns and find they align with
morphological boundaries the model was never given supervision for,
suggesting the regularizer recovers structure that is genuinely
present in the data rather than merely discarding capacity.
\end{abstract}

\begin{dedication}
To my parents, who never once asked when I would finish.
\end{dedication}

\chapter*{Acknowledgments}

I am deeply grateful to my advisor, Robert T. Halvorsen, for his
patience and for pushing me to ask sharper questions than the ones
I started with. Thanks also to my committee for their careful
reading of early drafts, and to the members of the Language
Processing Lab for years of useful arguments about what
``low-resource'' actually means. This work was supported in part by
a Dissertation Completion Fellowship from the Graduate College.
Finally, thanks to my family and friends for their patience during
the long stretches when this thesis was the only thing I could talk
about.

\tableofcontents
\listoftables
\listoffigures

\chapter{List of Abbreviations}

\begin{symbollist*}
\item[NLP] Natural Language Processing.
\item[POS] Part of Speech.
\item[RNN] Recurrent Neural Network.
\end{symbollist*}

\mainmatter

\chapter{Introduction}
\label{ch:intro}

Sequence labelling underlies a wide range of natural language
processing tasks, from part-of-speech tagging to named-entity
recognition. Modern approaches rely on dense neural representations
trained end to end~\cite{vaswani2017}, an approach that performs
well given abundant labelled data but degrades sharply for the
majority of the world's languages, where such data does not exist.

This thesis asks a narrower question: can a sparsity-inducing
regularizer, applied to a modest recurrent architecture, close part
of that gap without requiring any additional labelled data or
cross-lingual supervision? We motivate the approach in
Chapter~\ref{ch:related}, describe the model and training procedure
in Chapter~\ref{ch:model}, and report experiments across five
languages in Chapter~\ref{ch:results}.

\chapter{Related Work}
\label{ch:related}

Structured sparsity has a long history in statistical learning,
predating its recent applications to neural networks by decades;
early work on adaptive systems~\cite{holland1975} already treated
sparse, structured representations as a route to generalization
under limited data. More recently, generative approaches to
representation learning~\cite{goodfellow2014} demonstrated that
architectural constraints, rather than raw data volume, can shape
what a model learns to represent. We build on both lines of work,
combining a structured sparsity penalty with a data-efficient
recurrent encoder.

\chapter{Model and Training Procedure}
\label{ch:model}

We use a two-layer bidirectional LSTM encoder over subword units,
with a group-lasso penalty applied to blocks of hidden units
corresponding to fixed-width character windows. The penalty
coefficient is tuned on a held-out set of 50 labelled sentences per
language, and all remaining labelled data is reserved for the final
evaluation reported in Chapter~\ref{ch:results}. Full hyperparameter
settings and training curves are given in
Appendix~\ref{app:hyperparams}.

\chapter{Experiments and Results}
\label{ch:results}

We evaluate on five typologically diverse languages selected for
limited labelled data availability. Across all five, the sparsity-
regularized model outperforms a dense baseline of matched capacity
by a wide margin when fewer than 1{,}000 labelled sentences are
available, with the gap narrowing as more data is added -- consistent
with the regularizer acting primarily as an inductive bias rather
than an increase in effective capacity.

\chapter{Conclusions}

The structured sparsity penalty proposed in this thesis
substantially improves sequence-labelling accuracy under severe data
scarcity, and the resulting sparsity patterns correlate with
morphological structure the model was never told about. Future work
should examine whether this correlation can be made explicit and
used to guide annotation effort toward the sentences a model finds
hardest to sparsify.

\appendix

\chapter{Hyperparameter Settings}
\label{app:hyperparams}

All models were trained with the Adam optimizer, a learning rate of
$3\times 10^{-4}$, and early stopping on validation loss with a
patience of ten epochs. The group-lasso coefficient was selected
from $\{10^{-4}, 10^{-3}, 10^{-2}\}$ per language via the held-out
tuning split described in Chapter~\ref{ch:model}.

\backmatter

\bibliographystyle{plain}
\bibliography{references}

\chapter{Vita}

Priya Anand received her B.S. in Computer Science from Purdue
University in 2018 and her M.S. from the University of Illinois at
Urbana-Champaign in 2021, where she continued on to complete her
Ph.D. in the Language Processing Lab.

\end{document}
